% answers_short.tex —— 2007 年·普通高等学校招生考试·江西卷·文科 · 简短答案速查
% 由 tools/make_short_answers.py 自 answers.tex 的 \gksolhead/\gkanswer 提取，
% 并与 2007/2007_数学真题_江西文_答案要点.md 逐题比对后生成，勿手改（改源 MD 后重跑）。

\begin{gknotice}
\textbf{使用说明}：本节只给最终答案便于快速核对，完整推导见解析版。答案为编写者独立推导，\textbf{非官方标准答案}。
\end{gknotice}

\gkpart{选择题}
\begin{center}\begin{tabular}{ccccccccccccc}
\toprule
题号 & 1 & 2 & 3 & 4 & 5 & 6 & 7 & 8 & 9 & 10 & 11 & 12 \\
\midrule
答案 & B & B & A & D & A & D & B & B & C & C & A & C \\
\bottomrule
\end{tabular}\end{center}

\gkpart{填空题}
\begin{itemize}[leftmargin=2.6em]
  \item[\textbf{13.}] $1$。
  \item[\textbf{14.}] $7$。
  \item[\textbf{15.}] $(1,4)$。
  \item[\textbf{16.}] A、B、C。
\end{itemize}

\gkpart{解答题}
\begin{itemize}[leftmargin=2.6em]
  \item[\textbf{17.}] （1）$c=\dfrac12$；（2）$\left(\dfrac{\sqrt2}{4},\dfrac58\right)$。
  \item[\textbf{18.}] （1）$\theta=\dfrac{\pi}{6}$，$\omega=2$；（2）$x_0=\dfrac{2\pi}{3}$ 或 $\dfrac{3\pi}{4}$。
  \item[\textbf{19.}] （1）$0.8$；（2）$0.492$。
  \item[\textbf{20.}] （1）证明见解析；（2）$\arctan\dfrac13$；（3）$\dfrac32$。
  \item[\textbf{21.}] （1）$8$；（2）$T_{2n}=\dfrac9{16}\left[1-\dfrac{8n+3}{3^{2n+1}}\right]$。
\end{itemize}
